% !TeX program = xelatex
\documentclass[UTF8,a4paper,zihao=-4,fontset=fandol]{ctexart}
\usepackage[margin=25mm,headheight=16pt]{geometry}
\usepackage{amsmath,amssymb,booktabs,longtable,tabularx,array,graphicx}
\usepackage[dvipsnames]{xcolor}
\usepackage{tikz}
\usetikzlibrary{arrows.meta,positioning,shapes.geometric,calc}
\usepackage{listings,caption,subcaption,fancyhdr}
\usepackage[hidelinks,unicode]{hyperref}
\usepackage{xurl}
\urlstyle{tt}
\setlength{\emergencystretch}{3em}
\setlength{\parskip}{.3em}
\renewcommand{\arraystretch}{1.28}
\newcolumntype{Y}{>{\raggedright\arraybackslash}X}
\pagestyle{fancy}\fancyhf{}
\fancyhead[L]{\small 多机协同控制实验报告}
\fancyhead[R]{\small 协作初稿}
\fancyfoot[C]{\thepage}
\captionsetup{font=small,labelfont=bf}
\lstset{basicstyle=\ttfamily\small,breaklines=true,columns=fullflexible,
 keepspaces=true,frame=single,rulecolor=\color{black!20},showstringspaces=false,
 backgroundcolor=\color{black!2},aboveskip=8pt,belowskip=8pt}
\newcommand{\code}[1]{\nolinkurl{#1}}
\newcommand{\pending}[1]{\textcolor{BrickRed}{\textbf{待补充：}#1}}
% 图片可晚于正文提交；缺图时明确标为占位，不伪装成现场证据。
\newcommand{\reportphoto}[3]{%
 \begin{figure}[htbp]\centering
 \IfFileExists{#1}{\includegraphics[width=.8\linewidth,height=.30\textheight,keepaspectratio]{#1}}%
 {\fbox{\parbox[c][32mm][c]{.82\linewidth}{\centering\small 现场图片待归档\\\code{#1}\\#3}}}
 \IfFileExists{#1}{\caption{#2}}{\caption{#2（原始图片待归档）}}\end{figure}}
\tikzset{block/.style={draw=black!60,rounded corners=2pt,fill=blue!4,
 align=center,font=\small,minimum height=10mm,text width=32mm},
 flow/.style={-{Stealth[length=2mm]},thick},
 note/.style={font=\small,align=center}}

% 所有成员协商后统一修改此文件。不得将示例姓名当作真实贡献。
\newcommand{\ReportClass}{待填写课程与班级}
\newcommand{\ReportTeam}{第一组（人形机器人与方式二）及协作小组}
\newcommand{\ReportDate}{2026年10月15日}
% 当前报告核对的代码基线；后续代码变更时同步更新。
\newcommand{\SourceRevision}{6b6db8c}

\begin{document}
\hypersetup{pageanchor=false}
\begin{titlepage}
\centering
\vspace*{22mm}
{\Large 智能机器人系统综合实验\par}
\vspace{12mm}
{\Huge\bfseries 基于视觉与地图导航的\par 多机协同控制\par}
\vspace{8mm}
{\Large 方案设计与实验报告（协作初稿）\par}
\vspace{18mm}
\begin{tabular}{rl}
课程/班级：&\ReportClass\\[6mm]
团队名称：&\ReportTeam\\[6mm]
成员姓名、学号：&见第\ref{sec:division}节分工表\\[6mm]
报告日期：&\ReportDate\\[6mm]
代码依据：&\texttt{\SourceRevision}
\end{tabular}
\vfill
\begin{minipage}{.88\linewidth}
\small 本报告汇总 TonyPi 抱物穿行与四类 TurtleBot2 跟随方案。本次重点撰写人形机器人、方式二建图避障跟随；其余方式由对应小组在独立源文件中补充。姓名、个人贡献与定量实验数据尚待各组核实。文中区分代码实现、离线验证、局部实机验证与完整比赛验收。
\end{minipage}
\end{titlepage}
\pagenumbering{Roman}
\hypersetup{pageanchor=true}
\section*{摘要}
\addcontentsline{toc}{section}{摘要}
本实验以一台 TonyPi 人形机器人和多台 TurtleBot2 为平台，研究抱物穿行、视觉航路控制及多方式跟随的系统集成。任务要求人形机器人抱持物体完成至少一次左转和一次右转，终点放置物体；各移动机器人以可区分的方式保持跟随，并满足避碰、目标保持和终点停靠要求。

人形机器人部分采用“动作组执行＋视觉反馈”的分层结构：以抓取末帧固定上肢姿态，复用下肢步态；先通过开环路线定位动作偏移问题，再实现灰色挡板道路居中，最新版本进一步采用 A/B/C 固定标记进行测距、对中与转角切换。方式二采用赛前 GMapping 建图、比赛 AMCL 定位、AR 观测变换至地图、历史轨迹回退生成目标以及 move\_base/DWA 导航，并设置独立速度门控。报告给出分工模板、接口约定、实现原理、已知验证范围和实验记录表，供各组并行补充。

现有仓库记录支持局部实机运行与部分软件验证，尚不足以证明整条路线及全队比赛验收完成。本文不将仿真结果、默认参数或设计目标作为实测成绩。

\noindent\textbf{关键词：}TonyPi；TurtleBot2；动作组；视觉标记；地图导航；多机协同

\tableofcontents
\clearpage\pagenumbering{arabic}
\section{实验任务与场地}
\subsection{任务要求与验收口径}
依据课程《第五--六章：基于视觉与 AI 定位的多机协同控制实验》\cite{course}，系统同时包含人形机器人搬运和多车跟随两项任务。TonyPi 从指定位置抱起目标物体，在障碍通道中至少完成两个转角（左右转各至少一次），到达终点后放下物体；TurtleBot2 保持跟随并在终点区域停靠。课程总体评价由现场竞赛与方案、源码报告构成，权重分别为80\%和20\%。手册中的任务平台数量写为3--4台 TurtleBot2，详细评分表按四台、四种方式设置；本报告为四种方式保留独立章节，实际参赛车辆与方式认定须由全队和教师确认。

\begin{table}[htbp]\centering\small
\caption{任务要求与对应记录}
\begin{tabularx}{\linewidth}{p{30mm}YY}\toprule
要求 & 系统对应内容 & 应提交的证据\\\midrule
抱持、转弯与放置 & 抓取动作、上肢锁定、A左转/B右转/C停止 & 连续视频、动作文件版本、决策日志\\
不接触障碍物 & 人形路线约束；TurtleBot代价地图与门控 & 接触次数、障碍位置、最小间距实测\\
保持跟随目标 & 标记身份、目标时效检查、失效处置 & 丢失次数、持续时间、重捕获记录\\
多种独立方式 & 四组分别说明观测来源和控制链 & 各组实现、接口及现场运行证据\\
协同到达终点 & 任务状态、各车停靠条件 & 全队同场运行记录与到达时间\\\bottomrule
\end{tabularx}\end{table}

\subsection{场地与感知约束}
现场使用灰色地板及相对连续的灰色挡板构成通道，路线包含直行、左转、直行、右转和终点区域。地板接缝、胶带及其他地面标识均视为无关背景，不能作为路线识别标准。目标物块由早期蓝色调整为红色；抱持后物块可能遮挡画面下部，因此抓取确认与前方导航需要分别配置观察区域或头部姿态。

黄色雪糕筒可作为早期道路居中方案的终点候选；最新标记路线以固定标记 C 的距离触发停止。两种终点条件属于不同程序配置，不应在同一次实验记录中混写。图\ref{fig:route}仅表达路径拓扑，方向和长度以现场测量为准。
\begin{figure}[htbp]\centering
\begin{tikzpicture}[x=1cm,y=1cm]
\draw[line width=10mm,gray!15] (0,0)--(4,0)--(4,2)--(8,2);
\draw[flow,blue!70!black] (0,0)--(4,0)--(4,2)--(8,2);
\foreach \p/\t in {(0,0)/起点,(4,0)/A：左转,(4,2)/B：右转,(8,2)/C：停止放置}
 {\fill \p circle(2pt);\node[above=5mm,font=\small] at \p {\t};}
\node[below,font=\small] at (2,0) {第一段直行};
\node[right,font=\small] at (4,1) {第二段};
\node[below,font=\small] at (6,2) {第三段直行};
\end{tikzpicture}

\caption{路线拓扑示意；不按比例绘制，地面标识不参与定位}\label{fig:route}
\end{figure}
\reportphoto{figures/site/corridor.jpg}{灰色挡板通道现场}{选用通道全景；注明拍摄位置、日期和行进方向。}

\section{团队组织与成员分工}\label{sec:division}
\subsection{小组之间的职责}
第一组负责本次撰写的人形机器人与方式二部分。其余采用“第$i$组”代称；下表为方便协作，暂将第二、三、四组分别对应方式一、三、四，需由对应小组确认后定稿。仓库部分历史 README 使用“本组负责方式三”的叙述，不能据此推断第一组成员身份。
\begin{longtable}{p{28mm}p{66mm}p{43mm}}
\caption{全队功能分工（人员信息待核实）}\\\toprule
功能组 & 工作范围与交付 & 协同接口/负责人\\\midrule\endfirsthead
\toprule 功能组 & 工作范围与交付 & 协同接口/负责人\\\midrule\endhead
第一组：人形机器人 & 动作录制与组合、抓取/放置、相机调试、视觉闭环路线、运行日志 & 目标背部标记、任务状态；姓名待填\\
第二组：方式一（暂定） & AR直接跟随实际代码、标签尺寸、速度控制、失去标记处置；核对与课程固定航路点方式的区别 & \code{follow1}；姓名待填\\
第一组：方式二 & 地图采集、AMCL、AR转地图、轨迹目标、move\_base与速度门控 & \code{follow2}；姓名待填\\
第三组：方式三（暂定） & RGB-D目标检测、身份关联、实时避障与独立验收 & \code{follow3}；姓名待填\\
第四组：方式四（暂定） & 前一台TurtleBot位姿/轨迹获取、队列间距与末车停靠 & \code{follow4}；姓名待填\\
联调与报告 & 网络/TF核对、各方式独立性确认、数据汇总、视频与报告合并 & 总负责人、各组对接人待填\\\bottomrule
\end{longtable}

\subsection{本小组内部职责登记}
目前未提供成员姓名及个人贡献，以下是可认领的工作项，\textbf{不是对实际个人贡献的认定}。一人可承担多项，一项可由多人共同完成；贡献比例仅在成员协商后填写。
\begin{longtable}{p{23mm}p{26mm}p{48mm}p{37mm}}
\caption{本小组成员与贡献登记}\\\toprule
姓名/学号 & 职责方向 & 具体工作及交付证据 & 贡献说明/确认\\\midrule\endfirsthead
\toprule 姓名/学号 & 职责方向 & 具体工作及交付证据 & 贡献说明/确认\\\midrule\endhead
待填写 & 动作与机体 & 录制抓取/放置，组合抱物步态，检查姿态衔接；对应动作文件与提交号 & 待成员核实\\
待填写 & 视觉与状态机 & 相机维护、标记检测、路线切换、窗口与日志；对应代码与运行记录 & 待成员核实\\
待填写 & 建图与导航 & Kinect/扫描调试、建图、AMCL、代价地图与DWA参数；地图及配置 & 待成员核实\\
待填写 & 跟随与协同 & AR地图目标、轨迹回退、速度门控、网络及任务状态；联调日志 & 待成员核实\\
待填写 & 测试与报告 & 数据采集、统计、图片归档、章节编写与排版；原始记录及文稿 & 待成员核实\\\bottomrule
\end{longtable}

\subsection{其他小组内部职责登记规范}
方式一、三、四的章节各预留一张成员分工表。各组填写“姓名、学号、具体任务、交付证据、合作边界”，避免只写“参与调试”或“负责报告”。最终由各组负责人核对，并在全队分工表登记对接人。

\section{系统架构与协同接口}
\subsection{分层控制}
人形机器人负责物体搬运与领行；TurtleBot负责各自的目标观测、定位和运动控制。不同机器人的执行器由各自进程独占，感知线程只提供观测，避免多个程序同时写舵机或底盘速度。图\ref{fig:architecture}描述方式二的完整数据链；任务心跳是可选适配接口，不代表所有人形入口已经发送心跳。
\begin{figure}[htbp]\centering
\begin{tikzpicture}[node distance=8mm and 9mm]
\node[block] (v) {TonyPi相机\\固定标记/道路观测};
\node[block,right=of v] (s) {路线状态机\\抓取、直行、转弯、放置};
\node[block,right=of s] (a) {动作组执行\\抱物上肢固定};
\draw[flow] (v)--(s);\draw[flow] (s)--(a);
\node[block,below=13mm of v] (m) {TurtleBot相机\\背部AR目标};
\node[block,right=of m] (t) {按时间戳转到map\\历史轨迹与回退点};
\node[block,right=of t] (n) {move\_base/DWA\\地图与实时扫描};
\draw[flow] (m)--(t);\draw[flow] (t)--(n);
\node[block,below=of n] (g) {许可与速度门控\\输出至底盘};
\draw[flow] (n)--(g);
\node[note,below=of m,text width=68mm,anchor=north west] (h) {任务状态须另行接入：\\通用competition支持心跳桥；\\最新marker\_route尚需适配};
\draw[flow,dashed] (h.east)-|(g.west);
\end{tikzpicture}

\caption{人形机器人与方式二跟随的数据链}\label{fig:architecture}
\end{figure}

\subsection{软件平台与版本边界}
TonyPi 使用 Python 3、OpenCV、NumPy、PyYAML及设备原装舵机SDK；不同树莓派镜像的串口、SDK和相机接口存在差异，应使用实际部署配置。TurtleBot课程虚拟机为 Ubuntu 16.04、ROS Kinetic、Python 2.7，功能包放置于 \code{~/catkin_ws/src}，不能直接用 Python 3 替换 ROS 节点解释器。

\begin{table}[htbp]\centering\small
\caption{关键接口及约束}
\begin{tabularx}{\linewidth}{p{37mm}p{40mm}Y}\toprule
接口 & 内容 & 使用约束\\\midrule
\code{/ar_pose_marker} & 带ID与时间戳的标记位姿 & ALVAR标记；尺寸、相机内参和坐标系须一致\\
\code{target_pose} & 目标在地图中的位姿 & 按观测时刻变换，保留原时间戳\\
\code{/scan} & 障碍扫描 & 实际可能来自Kinect深度转扫描；检查高度覆盖与时效\\
\code{map -> odom} & 全局定位变换 & 建图用GMapping；固定地图定位用AMCL，不能同时发布\\
\code{/competition/mission} & 搬运任务状态 & 需持续更新；开环和最新标记入口不能默认视为已接入\\
\code{/cmd_vel_mux/input/navi} & 方式二底盘导航速度 & 由独立门控输出；不与键盘遥控争用\\\bottomrule
\end{tabularx}\end{table}

\subsection{拓扑与跟随方式的区分}
“TonyPi到首车再到后车”的队列拓扑，与“标记/地图/RGB-D/前车轨迹”的跟随算法是两个维度。当前方式二默认直接观察 TonyPi 背部 AR 标记，再在地图中导航；也支持输入前车地图位姿。实验时应固定一个目标来源并记录，不能把两个适配器同时发布到同一个目标话题。

多机应使用唯一节点名、命名空间与TF名称，并确保时间与地图原点一致。虚拟机 NAT 地址可供宿主机 SSH 访问，不保证机器人能反向访问 ROS；联调需验证共同局域网、\code{ROS_MASTER_URI}与\code{ROS_IP}，报告不收录密码等登录凭据。

\section{第一组：TonyPi人形机器人抱物与视觉路线控制}
\subsection{方案演进与入口区分}
本部分的核心思路是保留人工录制且已经具备抱持效果的抓取动作，以动作组作为运动基本单元，将视觉观测作为动作选择与状态切换条件。机器人双足步态受负载、关节偏置及地面摩擦影响，固定动作次数容易累积航向误差，因此从开环演练逐步转向闭环控制。
\begin{table}[htbp]\centering\small
\caption{人形控制程序的用途与边界}
\begin{tabularx}{\linewidth}{p{36mm}YY}\toprule
入口 & 主要逻辑 & 本报告采用的地位\\\midrule
\code{open_loop.py} & stand、grab、定次数前进/左右转、release & 调试基线，不能补偿累计偏移\\
\code{competition.py} & 识别/抓取确认、固定或视觉路线、放置确认；支持人工stand到grab起步 & 综合反馈框架，需现场ROI与任务接口配置\\
\code{lane_trial.py} & 人工摆位后stand到grab，有限次数道路居中/纠偏 & 局部闭环试验，不包含完整两次路口任务\\
\code{maze_route/marker_route.py} & 抱起后依次靠近A、左转找B、右转找C、停止放置 & 最新独立路线实现；局部实机验证，完整路线待验收\\\bottomrule
\end{tabularx}\end{table}

\subsection{动作组格式与抱物步态组合}
\code{.d6a}文件实质是SQLite数据库，\code{ActionGroup}表按\code{Index}顺序保存动作帧。每帧由帧序号、持续时间（毫秒）及各\code{ServoN}位置刻度组成，历史文件存在16通道和18通道两种宽度。总线舵机刻度与头部PWM脉宽不是同一种单位，不应将“刻度500中位”套用到头部1200或1500的PWM参数。角度与位置的关系可写为
\begin{equation}
q_i=k_i(p_i-p_{0,i}),
\end{equation}
其中$p_i$为位置刻度，$p_{0,i}$为约定参考刻度，$k_i$包含方向和角度比例。刻度500只是常用中位参考，不能单凭此值认定机械关节角为零；本方案主要复用实录动作，不依赖完整关节绝对角标定。

令抓取动作最后一帧的上肢集合为$\mathcal A=\{6,7,8,14,15,16\}$，步态第$k$帧各舵机位置为$p_{k,i}^{\rm gait}$，则抱物组合动作定义为
\begin{equation}
p_{k,i}^{\rm carry}=\begin{cases}
p_{{\rm last},i}^{\rm grab},&i\in\mathcal A,\\
p_{k,i}^{\rm gait},&i\notin\mathcal A.
\end{cases}
\end{equation}
这一组合保持下肢帧序列、持续时间与其余通道不变，只替换六个手臂通道，减少转向步态破坏夹持姿态的可能性。历史版本分别组合过\code{go_forward}、\code{go_forward_fast}及快速转向动作；最新标记路线使用\code{grab_final_v2}和\code{catch_ball_*}，不能跨版本混用抓取末帧。

最新动作副本中的六个抱持刻度依次为566、908、40、435、108、935。\code{mkcopies.py}从设备原始动作生成独立副本，并逐帧核对未修改的腿部、索引和时间。放置动作仅在第一帧下蹲时保持抱紧手臂，后续降下、松开及站起帧保留原动作。运行时若首帧与当前姿态差值超过25刻度，将过渡时长延长到300\,ms；这一过渡不等于动态稳定性已获证明。

针对机体前倾，仓库另存两组大腿舵机相反方向的20刻度偏置候选，均带\code{_offset}标记。该方案未作为当前机器人的默认步态，也没有完整实机稳定性结论，不能在报告中写成最终校正量。

\subsection{抓取、持物观察与头部姿态}
早期通用视觉框架采用HSV/LAB分割、形态学去噪及轮廓筛选；物块改为红色后使用跨越HSV色相两端的双阈值。通用框架可在抓取后检查持物ROI，失败时经复位再识别，并在放置后检查手中区域与地面区域。人工起步模式则从stand直接衔接grab，省去搜索与接近阶段。

最新\code{marker_route}默认原地播放\code{grab_final_v2}后锁定上肢，\textbf{并未接入上述颜色持物确认、自动失败重抓和放置结果确认链}。因此“抓取动作完成”与“物块已可靠抓住”应分别记录，不能将通用框架的功能当作最新入口已验证的能力。比赛要求掉落后回到掉落点重新抱起；当前路线程序也没有完成该恢复闭环。

头部参数按程序和目标分别标定：早期道路/持物调试使用PWM servo1=1200，servo2保持当前值；最新标记路线默认pitch=1500、pan=1482，用来观察前方墙面标记。不同俯仰角会改变地面可见区域、物块遮挡与近距离码出画面的位置，报告应为每次运行登记真实参数，不能把两套姿态并写成同一配置。

\subsection{早期灰色挡板道路居中}
地面与挡板都为灰色，仅靠颜色阈值难以区分。早期道路算法使用边缘、线段、挡板支撑立柱与底边的几何关系，排除孤立地缝和地面胶带。抱持物块所在区域单独遮罩，避免把红色物块边缘误判为道路边界。

在近、远两个图像截面分别得到左右边界$x_L(y),x_R(y)$，可定义归一化横向误差与方向误差
\begin{align}
e_c&=\frac{x_L(y_n)+x_R(y_n)}{2W}-\frac12,\\
e_h&=\frac{x_L(y_f)+x_R(y_f)-x_L(y_n)-x_R(y_n)}{2W}.
\end{align}
其中$W$为图像宽度。组合误差超过阈值时选择左/右纠偏动作，误差处于容差内时执行一次前进，等待动作结束后重新观察。观测需连续稳定，不能在一张模糊帧上立即切换。挡板同色、遮挡、低对比和双足运动模糊仍可能导致无效道路判断；该方法不提供米制障碍距离或通用绕障规划。

\subsection{最新A/B/C固定标记识别与测距}
最新程序不依赖地面标识，而是在墙面布置三张已知模板标记。每张标记由$9\times9$格构成，外两格为黑边，中间$5\times5$格为数据，黑色正方形边长14\,cm。检测过程为自适应阈值、四边形轮廓筛选、透视展开、单元格采样，以及四个旋转方向的模板匹配；允许的数据位误差为1格。它们是程序自定义模板，\textbf{不能当作ArUco、AprilTag或TurtleBot所用ALVAR标签互换}。

以标记四角为三维已知点，结合相机内参$K$与畸变参数，通过PnP求取旋转$R$和平移$t$：
\begin{equation}
\lambda\begin{bmatrix}u\\v\\1\end{bmatrix}
=K\left[R\mid t\right]\begin{bmatrix}X\\Y\\Z\\1\end{bmatrix}.
\end{equation}
检测模块输出摄像头到标记中心的直线距离、方位角及墙面朝向。内参读取设备上的相机标定文件。码中心与镜头存在高度差时，PnP直线距离不能直接等同于沿地面水平距离，转角距离需用真实摆位复核。

\subsection{完整路线状态机与自适应纠偏}
路线状态依次为准备、抱起、走向A、左转寻找B、走向B、右转寻找C、走向C、停止与放置。距离阈值默认40\,cm；路口并非固定播放若干次就认定90度转弯，而是沿预定方向逐次转动，直到下一张目标码进入视野，再由直行阶段继续对中。每次转前后分别等待300\,ms，每次前进后等待200\,ms。

为降低步态大转角引起的左右摆动，程序对连续3帧观测取中位数，并根据近期单次转向后的方位变化估计转角$\hat\theta_{\rm step}$。纠偏门限采用
\begin{equation}
|\beta|>\max\!\left(8^\circ,0.6\hat\theta_{\rm step}\right),
\end{equation}
其中$\beta$为目标方位角。程序还从可见标记时的相邻步观测估计直行偏航，并在目标暂不可见时进行预测补偿。\textbf{标记不可见与相机无新帧是不同故障}：前者在计数约束下仍可能继续走，后者等待恢复。默认未见目标的盲行上限为25步，目标曾可见后连续丢失也有暂停阈值；这一实现与课程“丢失后原地重捕获”的严格要求存在差距，需在正式联调前收紧失效策略并验收。

当标记在近目标区域居中消失时，程序可使用最近观测判定已到达。这是处理近距离码从画面下缘离开的启发式，不能替代障碍测距。路口最多尝试12次转向；连续20步推进不足5\,cm、总步数过多或相机长时间无图像时触发暂停/退出。默认没有深度或测距输入，因此尚不能声称具备完整防碰撞保证。

\subsection{相机可靠性与可视化调试}
现场曾出现\code{/dev/video0}消失，内核提示USB枚举错误和线缆异常，重新插接后恢复。早期通用框架通过本地HTTP快照共享网页相机，避免重复打开已被占用的设备；最新\code{robust_camera.py}则按USB身份查找相机、处理网页相机释放和断线重连，只返回0.5秒内的新画面。两种相机管理模式应各自独立运行。

最新窗口显示运行阶段、目标标记框、标记ID、距离与方位等信息；支持S启动/恢复、P暂停、Q或Esc退出。每次运行保存\code{route.log}、\code{decisions.csv}及标注截图，用于核对状态转移的观测依据。截图中识别框和数值只证明该时刻的检测，不证明全程搬运成功。
\reportphoto{figures/humanoid/marker_debug.jpg}{标记航路的实机窗口}{选用已提供的屏幕照片原件，保留标记框和日志；避免以照片读数替代日志统计。}

\subsection{已有验证与后续验收}
最新仓库说明记录实机已跑通“抱起、走向A、在约39.6\,cm处开始左转”；完整A/B/C路线仍未实机验收\cite{maze}。仓库测试使用合成透视图验证标记识别与PnP，并用给定步长、偏航和转角模型演练路线；合成实验精度与路线完成不能作为实地成功率。

后续应重点记录：满载单步距离与偏航、左右转角分布、两次路口完成率、物块掉落/接触次数、相机断线频次及完整路线时长。抓取成功确认、丢码原地重捕获、掉落恢复、终点放置确认以及任务心跳适配，仍需在最新标记入口中形成一致闭环。

\section{第二组（暂定）：方式一AR直接跟随}
\subsection{已有实现摘要与方式认定}
仓库已合并从虚拟机提取的\code{ar_follower}包及独立实验说明\cite{follow1}。该实现通过Kinect检测前方机器人背部ALVAR标记，取标记光学坐标中的前向距离$z$和横向偏差$x$，以比例控制输出速度：
\begin{equation}
v=\operatorname{sat}\!\left(K_v(z-d_0)\right),\qquad
\omega=\operatorname{sat}\!\left(-K_\omega x\right).
\end{equation}
实际符号以节点坐标约定为准。默认目标ID为1、跟随距离0.4\,m、最大线速度0.3\,m/s、最大角速度1.0\,rad/s。前方60度扇区障碍小于0.25\,m时停止前进并慢速左转。该节点直接输出到\code{/cmd_vel_mux/input/teleop}，不经过方式二的速度门控。

课程参考的方式一强调固定场地标签航路点，而该组已提交程序属于移动目标背贴AR直接跟随，两者存在差别。应由对应组明确最终参赛定义、与方式二/三的独立性及教师确认情况，不在本初稿中强行等同。

\subsection{本组应补充的内容}
\pending{填写实际组号、现场标签尺寸、距离标定、完整跟随记录、转角与遮挡试验。现有节点缺少扫描/目标消息超时看门狗，不能将收到空检测消息时停车等同于所有通信中断时停车。}
\begin{table}[htbp]\centering\small
\begin{tabularx}{\linewidth}{p{28mm}YY}\toprule
姓名/学号 & 具体任务 & 交付证据/合作说明\\\midrule
待填写 & 标记与相机标定 & 待填写\\
待填写 & 跟随控制与实机测试 & 待填写\\
待填写 & 数据整理与报告 & 待填写\\\bottomrule
\end{tabularx}\caption{第二组（暂定）内部职责}\end{table}

\section{方式二：建图定位与避障跟随}
\subsection{建图避障跟随的实现总览}
\subsubsection{方案与功能包组织}
方式二将“看见目标”与“安全到达跟随位置”分开：相机观测TonyPi背部AR标记，在地图中恢复其位置并积累历史轨迹；随后由导航栈在已知地图和实时障碍约束下生成底盘速度。目标朝向不直接作为跟随方向，避免背贴标签朝向与实际行进方向不一致时生成错误目标。

独立ROS功能包\code{multi_robot_following}位于\code{turtlebot/follow2/}，部署到虚拟机\code{~/catkin_ws/src}。包内包括建图入口、AMCL/move\_base配置、目标源适配、轨迹跟随、速度门控与任务心跳桥，不覆盖原有课程包。

\subsubsection{沿用课程建图链与地图质量检查}
底盘启动与键盘遥控已由小组确认可用：\code{turtlebot_bringup minimal.launch}及\code{turtlebot_teleop keyboard_teleop.launch}。建图优先沿用课程\code{turtlebot_navigation gmapping_demo.launch}，另用课程导航RViz配置查看地图、扫描和TF。该入口已集成部分相机和导航节点，不与新包的建图入口重复启动。

建图时采集里程计和扫描，通过GMapping估计二维栅格地图；完成后使用\code{map_saver}成对保存\code{.yaml}与\code{.pgm}。比赛导航加载固定地图，AMCL估计当前车位姿，move\_base进行规划。场地障碍布局改变需重新建图；从建图切换到定位时，必须停止GMapping，避免与AMCL竞争\code{map -> odom}。

现场曾报告约35\,cm高障碍未在图中稳定出现。不能只按URDF假定镜头高度便认定扫描必然覆盖：应分别检查原始深度图、深度转扫描的采样行与范围、实际安装俯仰、TF、有效近距，以及RViz中\code{/scan}是否已有对应点。若扫描本身漏检，调节地图膨胀层不能补回未观测障碍；若扫描有点而地图无障碍，再排查建图、TF和清除配置。当前尚未提供漏检已解决的完整实测记录。
视频8\,s处可见多台搭载电脑的 TurtleBot2 和现场标签布置，可用于说明实际平台外观。
\reportphoto{figures/follow2/robot.jpg}{视频8\,s处的 TurtleBot2 平台与现场布置（V008）}{原始视频截图。}
图中车辆与方式编号尚无法对应，故本图只作为平台实物记录，不把其中某台车认定为方式二实车，也不据画面推算镜头高度或标记尺寸。

\subsubsection{AR目标观测到地图坐标}
方式二默认\code{target_source:=marker}，选取指定ID的ALVAR标记。\code{marker_target.py}验证时间戳、有限数值与量测范围，在观测时刻$t_m$将目标变换到地图：
\begin{equation}
{}^{m}T_{\rm target}(t_m)
={}^mT_c(t_m)\,{}^cT_{\rm target}(t_m).
\end{equation}
其中$c$为相机观测坐标系，$m$为地图坐标系。输出保留原始量测时间，TF缺失、过期或外推失败时丢弃观测，不以当前时间重新包装旧数据。节点默认量测时效0.5\,s、最大量测范围4\,m，这些是代码限制，不是实际检测精度。

另一适配器可订阅前车地图定位位姿，但必须共享地图原点并检查时间戳、跳变和协方差。当前第一组方案以直接观察TonyPi为主；只有确认切换到前车定位模式时，才在实验表中登记该数据源。

\subsubsection{历史轨迹与沿路径回退的跟随点}
\code{Trail}将有效目标位置记录为有界折线$P=(p_0,\dots,p_n)$。默认最小采样间距0.08\,m、最大保留路径长度20\,m、单次跳变上限0.8\,m。跟随距离默认$d_f=0.9$\,m，从最新点沿折线反向累计弧长；若回退点位于$p_{j-1}$与$p_j$之间，令该段长度$\ell_j=\|p_j-p_{j-1}\|$、剩余距离$r$，则
\begin{equation}
p_g=p_j+\frac{r}{\ell_j}(p_{j-1}-p_j).
\end{equation}
轨迹不足$d_f$时不生成有效目标。与沿目标当前朝向直接平移相比，这种方法保留已观测路径的转角信息。但move\_base仍会在可通行区域平滑或重规划，不能宣称精确重放前车全部轨迹，也不能保证后车永远不抄近路。
\begin{figure}[H]\centering
\begin{tikzpicture}[x=1cm,y=1cm]
\draw[gray!45,line width=7mm] (0,0)--(3,0)--(3,1.6)--(7,1.6);
\draw[blue!65!black,thick] (0,0)--(3,0)--(3,1.6)--(7,1.6);
\foreach \x/\y in {0/0,1/0,2/0,3/0,3/.8,3/1.6,4/1.6,5/1.6,6/1.6,7/1.6}
 \fill[blue!65!black] (\x,\y) circle(1.7pt);
\fill[red!70!black] (7,1.6) circle(3pt);
\node[above,font=\small] at (7,1.6) {目标$p_n$};
\fill[green!50!black] (4,1.6) circle(3pt);
\node[above,font=\small] at (4,1.6) {回退点$p_g$};
\draw[flow,red!60!black] (6.8,2.3)--(4.2,2.3);
\node[above,font=\small] at (5.5,2.3) {沿折线回退$d_f$};
\draw[fill=orange!30] (1,0) circle(3mm);
\node[below,font=\small] at (1,0) {当前后车};
\node[below,font=\small] at (5,.6) {最终行驶路径仍由导航栈规划};
\end{tikzpicture}

\caption{沿目标历史折线回退生成导航目标；示意图不含真实场地尺度}
\end{figure}

历史折线回退沿目标已经走过的路径生成目标点，转角处不会简单沿当前视线把机器人拉向前车。图中回退距离由累计弧长定义，区别于两点直线间距。得到回退点后仍需执行足迹与地图检查，路径几何本身不保证底盘能够安全通过。


\subsubsection{代价地图、DWA与目标可行性}
全局规划依托固定地图及目标点；全局与局部代价地图接收实时扫描，DWA在速度与加速度约束下评估短时轨迹。候选目标在发送前执行圆形足迹栅格检查，越界、未知或高占用栅格均不可进入。足迹、膨胀半径及障碍更新参数必须按真实车体和通道标定，默认值不是经过测量的保证。
\begin{table}[H]\centering\small
\caption{方式二部分默认参数（非最终实测值）}
\begin{tabularx}{\linewidth}{p{52mm}p{22mm}Y}\toprule
参数 & 默认值 & 作用\\\midrule
跟随路径间距 & 0.9\,m & 与目标沿历史路径保持滞后\\
最大前进速度 & 0.15\,m/s & 限制跟随底盘速度\\
最大角速度 & 0.5\,rad/s & 限制转向速度\\
线加速度限制 & 0.3\,m/s$^2$ & DWA动态约束\\
控制/规划频率 & 10/1\,Hz & 局部控制与全局重规划\\
DWA预测时长 & 1.5\,s & 轨迹评估窗口\\
平面目标容差 & 0.15\,m & 到达导航目标的容差\\\bottomrule
\end{tabularx}\end{table}
当前配置禁止倒车，并关闭自动原地旋转恢复，以降低拥挤通道中盲目恢复的风险。其代价是某些受阻情形需要重新规划或人工复位，而不能自行摆脱。视野之外、深度缺测或低于有效扫描面的物体仍可能漏检，导航栈不等于全空间避碰。

\subsubsection{双层许可与速度门控}
跟随控制层检查目标、定位、TF、地图/代价地图、任务状态与目标最小间距，满足条件才发导航目标与运动许可。独立\code{velocity_gate.py}进一步要求许可、导航速度和扫描均新鲜，检查前方扇区有效量测并拒绝倒车、NaN/Inf及非法速度。可概括为
\begin{equation}
u_{\rm out}=\begin{cases}
\operatorname{limit}(u_{\rm nav}),&G_{\rm mission}\land G_{\rm target}\land
G_{\rm loc}\land G_{\rm map}\land G_{\rm TF}\land G_{\rm scan}\land G_{\rm cmd},\\
0,&\text{其他情况}.
\end{cases}
\end{equation}
其中$G$表示相应许可；它们分别由上游与独立门控实现，并非一个节点完成所有检查。NaN/Inf占满的扫描扇区被视为缺测，不当作空闲空间。门控具有壁钟检查线程，以避免ROS仿真时钟停顿后速度无法及时归零。

AMCL在机器人静止时不一定持续发布位姿，程序定期请求无运动定位更新，并配置重采样频率，使新扫描可更新定位结果。请求本身不制造有效扫描或虚假新鲜度，定位故障仍需在RViz重新初始化或排查传感器。

\subsubsection{任务状态联动与当前集成缺口}
跟随器需要持续收到允许阶段的\code{/competition/mission}。通用\code{competition.py}可通过UDP心跳桥提供状态；\code{open_loop.py}不发送该心跳，最新\code{marker_route.py}同样未实现该接口。因此“两个程序分别能启动”不等于“已经自动协同”：需增加适配器，将标记路线的搬运、放置、完成、暂停和故障映射到任务状态，明确频率、超时和停止行为后再联调。

单车导航调试时可以在终端持续发布\code{carrying}，但这仅是测试输入，不能写成全队自动联动已完成。终点还需核对后车是否真正进入终点区域并停稳，不能只用TonyPi的放置完成判定所有车辆到达。

\subsubsection{验证范围与待测指标}
已有部署记录表明功能包在ROS Kinetic/Python 2.7下完成构建、语法与launch解析检查；隔离ROS master中，在缺少目标、定位、任务与扫描时观测到零速度和禁止运动。该检查未向真实底盘发送导航速度，证明的是部分失效路径的逻辑，不能推导实际绕障成功率。

本组已确认底盘和键盘遥控可用，地图质量、AMCL稳定性、动态避障、目标丢失后原地恢复及全程跟随仍需按第\ref{sec:experiments}节采集真实数据。尤其应在35\,cm障碍、转角目标遮挡和网络中断条件下逐项复核，而非直接进行多车混合测试。

\section{第三组（暂定）：方式三RGB-D实时检测与避障跟随}
\subsection{已有实现摘要}
仓库\code{turtlebot/follow3/rgbd_following}已包含颜色与深度候选检测、多帧目标关联、外观模板及有界重识别、历史轨迹目标、滚动里程计代价地图与独立速度门控\cite{follow3}。其观测来源为实时RGB-D与本车里程计，不依赖预建地图或前车广播位置，区别于方式二的固定地图导航。

现有说明将检测阈值作为示例，要求登记真实身份模板并在独立样本上评估。区域深度有效不等于识别到了TonyPi；颜色候选也不能替代身份连续性验证。源码和合成演练已提供，真实目标识别、导航与全队终点验收仍需对应组确认。
\subsection{本组补充区域}
\pending{实际组号、识别特征与数据集、RGB/深度配准、遮挡重识别策略、主动绕障方式、默认参数与实测结果。应给出误识别、目标切换、目标丢失及终点停车的原始证据。}
\begin{table}[htbp]\centering\small
\begin{tabularx}{\linewidth}{p{28mm}YY}\toprule
姓名/学号 & 具体任务 & 交付证据/合作说明\\\midrule
待填写 & 数据与目标识别 & 待填写\\
待填写 & 控制、避障与联调 & 待填写\\
待填写 & 评估与报告 & 待填写\\\bottomrule
\end{tabularx}\caption{第三组（暂定）内部职责}\end{table}

\section{第四组（暂定）：方式四前车轨迹与队列跟随}
\subsection{方案接口}
方式四以队列中的前一台TurtleBot为目标，期望保持固定间距或队形。仓库\code{turtlebot/follow4/README.md}给出复用方式二地图位姿输入的接入约定：前车发布唯一定位话题，各车共享地图并配置独立命名空间/TF；也允许由队友明确实现背贴标记观测。队列拓扑应指定每台车的前车，避免末车错误跟随其他运动物体。

现有目录主要为接入说明，不足以证明一套独立方式四控制器已完成实机运行。若复用方式二，应由全队明确算法差异和计分认定，不能仅因目标从TonyPi换成TurtleBot就自动宣称四种方式独立。
\subsection{本组补充区域}
\pending{实际组号、前车身份、目标输入与频率、轨迹保持/队形算法、失联处置、队列长度对误差的影响，以及末车终点停靠结果。}
\begin{table}[htbp]\centering\small
\begin{tabularx}{\linewidth}{p{28mm}YY}\toprule
姓名/学号 & 具体任务 & 交付证据/合作说明\\\midrule
待填写 & 前车位姿/身份接口 & 待填写\\
待填写 & 队列跟随与避碰 & 待填写\\
待填写 & 多车测试与报告 & 待填写\\\bottomrule
\end{tabularx}\caption{第四组（暂定）内部职责}\end{table}

\section{实验设计、证据与结果记录}\label{sec:experiments}
\subsection{验证层级与当前事实}
本报告不将“已编写”“离线通过”“局部实机”和“完整验收”混为一谈。测试数量随仓库更新会变化，自动化测试通过也不能证明实际无碰撞。正式结果须绑定代码提交、设备、参数、场地与原始证据。
\begin{longtable}{p{32mm}p{54mm}p{51mm}}
\caption{已有证据及其范围}\\\toprule
项目 & 当前证据 & 能说明/不能说明\\\midrule\endfirsthead
\toprule 项目 & 当前证据 & 能说明/不能说明\\\midrule\endhead
动作副本 & 数据库帧、时间与非手臂通道核对；独立副本生成脚本 & 说明组合遵循指定规则；不证明抱持稳定性\\
人形局部路线 & 最新README记录抱起、走到A并约39.6\,cm处左转 & 支持首段局部运行；不证明B/C路线完成\\
人形离线验证 & 合成标记透视图与给定步态模型路线演练 & 说明设定条件下算法逻辑；不是真实误差分布\\
相机链 & 历史恢复检查连续取得15帧640\,$\times$\,480新图像 & 说明当次取流恢复；不保证持续运动不掉线\\
TurtleBot基础 & 小组确认底盘启动与键盘遥控可用 & 说明基础驱动可用；不等于自主导航成功\\
方式二隔离检查 & Kinetic构建、launch解析、缺少必要输入时零速度 & 部分软件与停止逻辑；未完成实地绕障验收\\
完整比赛 & 未提供完整连续运行及统计记录 & 成功率、总时长、碰撞与掉落次数均待测\\\bottomrule
\end{longtable}

\subsection{分阶段实验与异常注入}
\begin{enumerate}
\item \textbf{单动作与视觉。}人工监督下分别测试站立、抓取、抱物前进、左右转和放置；固定相机姿态后测量标记距离/方位，统计多次误差与失效距离。
\item \textbf{人形单段与完整路线。}先测直行偏移，再测A左转/B右转，最后无人工干预完成A/B/C全程。分别记录动作次数、转角偏差、接触、掉落及相机断线。
\item \textbf{地图和单车导航。}确认35\,cm障碍进入原始深度与扫描，再检查地图。先静态目标导航、定位丢失和规划失败，再接入移动AR目标。
\item \textbf{目标与任务失效。}遮挡标签、停止目标消息、停止任务心跳、断开扫描或使TF过期，逐项测量停止延迟。人形最新程序的盲行与近距丢码启发式应单独评估。
\item \textbf{多车联调。}固定队列拓扑、共享地图与时间基准，避免速度控制源冲突；记录各车到达终点条件，不只记录领行机器人结束。
\end{enumerate}

\subsection{指标定义}
距离标定用均方根误差$\operatorname{RMSE}_d=\sqrt{N^{-1}\sum_i(\hat d_i-d_i)^2}$；$d_i$必须是与PnP输出同一定义的真实距离。跟随间距误差定义为$e_f(t)=d(t)-d_0$，报告平均绝对误差与最大偏差，同时说明$d$采用欧氏距离还是路径弧长。完整任务成功率定义为$N_{\rm success}/N_{\rm trials}$，每次试验均需记录，不能剔除失败试跑而不说明。

对接触次数、掉落次数、目标丢失持续时间、最长无新帧时间、停止延迟和各车到达时间分别统计。转角是否达到路线要求应通过外部测量或可复核视频判定，不能把“看见下一张码”直接记为精确90度。

\begin{table}[htbp]\centering\small
\caption{正式结果汇总模板；破折号表示未测，不表示零}
\begin{tabularx}{\linewidth}{Yp{26mm}p{26mm}p{25mm}}\toprule
项目 & 样本数 & 结果 & 证据编号\\\midrule
标记距离误差 & 待测 & --- & 待填写\\
满载单步距离/偏航 & 待测 & --- & 待填写\\
人形完整路线成功率 & 待测 & --- & 待填写\\
方式二完整跟随成功率 & 待测 & --- & 待填写\\
35\,cm障碍检测与避让 & 待测 & --- & 待填写\\
目标/任务失效停止延迟 & 待测 & --- & 待填写\\
全队到达率/总时长 & 待测 & --- & 待填写\\\bottomrule
\end{tabularx}\end{table}

\subsection{数据归档与照片使用}
运行记录模板保存在\code{latex/data/experiment_records.csv}。每行对应一次试验，保存代码版本、设备、运行参数、结果及证据路径；原始视频和bag较大时放在团队共享存储中，报告记录可访问的证据位置。不要将不同机器人的动作参数或不同版本的相机设置合并统计。

用户提供的通道、实机窗口及TurtleBot照片可作为场地和调试过程说明；应将原图归档到\code{figures/site}、\code{figures/humanoid}和\code{figures/follow2}并注明出处。当前源码支持图片尚未归档时显示明确占位框；占位框不是实验结果。终稿应补齐实际引用照片或移除相应图位。

\section{问题分析与总结}
第一组围绕人形抱物和地图避障跟随建立了可分步调试的实现。动作层通过固定抓取末帧上肢、复用原步态下肢实现抱物前进与转向；视觉层从开环次数控制发展到道路几何居中及固定标记航路，减轻直行偏航的累计影响。方式二将AR相对观测变换到地图，利用历史路径回退生成目标，再交由导航栈规划，并以任务、定位和扫描等有效性条件约束速度输出。

实验过程表明，影响系统的因素不仅是控制算法，还包括动作版本、抱持负载、相机姿态、USB供电与连接、深度扫描覆盖、多机TF与网络。仅根据模型镜头高度推断障碍可见、仅根据动作次数认定转角或仅根据单帧识别认定完整跟随成功，都会使报告结论超出证据。

当前主要未完成项为最新标记路线的完整实机验收、抓取/掉落/放置反馈的一致接入、严格的丢目标原地恢复、任务心跳适配、矮障碍检测复核及全队同场结果统计。后续应以日志和连续视频建立可复核记录，再填写性能结果和成员最终贡献。其余小组应在各自章节补足实际实现与方式独立性说明，共同完成终稿。

\begin{thebibliography}{9}
\addcontentsline{toc}{section}{参考资料与实现依据}
\bibitem{course} 课程实验手册．第五--六章：基于视觉与AI定位的多机协同控制实验．仓库根目录同名\code{.docx}文件。
\bibitem{repo} 项目协作仓库．\url{https://github.com/yky666/intelligent_robot_system}．本次代码核对基线：\code{6b6db8c}，后续修改见版本历史。
\bibitem{maze} TonyPi标记航路实现及实验状态．\code{humanoid/maze_route/README.md}，\code{marker_route.py}，\code{markers.py}，\code{robust_camera.py}，\code{tests/test_maze_route.py}。
\bibitem{follow2} 第一组地图跟随实现与部署记录．\code{turtlebot/follow2/multi_robot_following/}，\code{docs/vm_follow2.md}，\code{docs/turtlebot_debug_commands.md}。
\bibitem{follow1} AR直接跟随实验总结与提取功能包．\code{docs/ar_marker_follow_experiment_summary.md}，\code{turtlebot/follow1/ar_follower/}。
\bibitem{follow3} RGB-D实时跟随方案、状态与验收要求．\code{turtlebot/follow3/README.md}及对应功能包。
\end{thebibliography}

\appendix
\section{复现入口与协作清单}
\subsection{源码入口索引}
\begin{longtable}{p{31mm}p{110mm}}
\toprule 内容 & 仓库路径\\\midrule\endhead
最新人形路线 & \code{humanoid/maze_route/marker_route.py}\\
检测与相机 & \code{humanoid/maze_route/markers.py}，\code{robust_camera.py}\\
动作副本 & \code{humanoid/maze_route/actions/}；生成规则见同目录\code{mkcopies.py}\\
通用反馈与道路 & \code{humanoid/scripts/competition.py}，\code{lane_control.py}，\code{lane_trial.py}\\
方式二入口 & \code{turtlebot/follow2/multi_robot_following/launch/vm_follow2.launch}\\
目标与导航控制 & 该包\code{scripts/marker_target.py}与\code{map_follower.py}\\
AR控制与共享门控 & \code{turtlebot/follow1/ar_marker_following/scripts/marker_follower.py}与\code{velocity_gate.py}\\
轨迹与几何核心 & 新独立包\code{turtlebot/follow2/competition_core/src/planning.py}；旧包仍保留\code{src/competition_core/planning.py}，安装布局待核对\\
导航默认参数 & 该包\code{config/move_base.yaml}及\code{costmap_common.yaml}\\\bottomrule
\end{longtable}

\subsection{只查看最新标记检测}
在具备相机、标定文件及显示环境的TonyPi上，从完整仓库运行：
\begin{lstlisting}[language=bash]
cd ~/intelligent_robot_system
python3 humanoid/maze_route/marker_route.py --no-servos
\end{lstlisting}
\code{--no-servos}不访问舵机串口并隐含dry-run；单独\code{--dry-run}仍可能设置头部，二者不能混淆。真实运动入口会关闭厂家上位机并停止相关服务，运行前核对同目录README与设备配置；本报告构建不会运行机器人控制脚本。

\subsection{方式二现场启动索引}
底盘、相机、扫描与地图请按\code{docs/turtlebot_debug_commands.md}逐终端启动，避免重复驱动。最新结构还需同时部署方式一共享包和独立核心包，并解决正文记录的安装问题。完成迁移后构建、地图及AMCL初始位姿检查，主要跟随入口为：
\begin{lstlisting}[language=bash]
source /opt/ros/kinetic/setup.bash
source ~/catkin_ws/devel/setup.bash
roslaunch multi_robot_following vm_follow2.launch \
  map_file:=$HOME/catkin_ws/maps/competition.yaml \
  marker_id:=1 marker_size_cm:=9.0
\end{lstlisting}
标记尺寸9.0\,cm是该launch默认值，必须按实际ALVAR黑边尺寸替换；与人形航路14\,cm自定义标记无关。正式跟随还需要持续任务状态和各项许可，单独启动此命令不会自动形成全队协同。

上述是入口索引，不是最新拆包后的实机成功记录；本次未在ROS Kinetic设备重新部署或运行这些命令。

\subsection{报告提交前核对}
\begin{enumerate}
\item 各组确认组号、成员姓名学号及真实个人贡献，修改分工表和封面信息。
\item 补齐照片原件、地图/参数版本、每次试验记录与连续运行证据。
\item 将合成测试与实测结果分开，填入样本数、失败次数及实际验收结论。
\item 核对四种方式的独立性、任务状态适配及终点停靠定义。
\item 编译并检查PDF目录、公式、表格和图片；删除或补齐所有待填写项目。
\end{enumerate}

\end{document}
